Pour apprécier la contribution possible des des membres des associations IHEDN, le comité a choisi de raisonner selon les trois phases du cycle de crise. Une telle approche permet de distinguer ce qui relève de la préparation, ce qui peut être envisagé pendant l'événement, et ce qui prend sens dans l'analyse \textit{a posteriori}.

\subsection{Avant la crise : préparation et résilience}

En amont, les membres de l'\ARSeize{} peuvent contribuer à diffuser une culture de sécurité, à sensibiliser les populations et à participer à des dispositifs de préparation, notamment autour des \gls{pcs}\footcite{PCS_PICS_def} ou d'exercices à destination des populations\footcite{Gestion_de_crise_def_exemples}. Cette contribution relève moins d'une intervention de crise que d'un travail de préparation, de mise en relation et de partage d'expérience.


\begin{proposition}{Développer une culture de préparation et de résilience}
L'\ARSeize{} pourrait contribuer à la diffusion d'une culture de préparation aux crises en s'impliquant dans des actions de sensibilisation, de formation et d'information à destination des élus, des auditeurs et du grand public, notamment autour des dispositifs de gestion de crise et des enjeux informationnels contemporains.
\end{proposition}

Le tableau \ref{tab:sensibilisation_resilience} présente six axes d'action pouvant être envisagés afin de favoriser les échanges autour des enjeux et problématiques liés à la gestion de crise.

\begin{table}[H]
\centering
\footnotesize
\setlength{\tabcolsep}{4pt}

\resizebox{\linewidth}{!}{%
\begin{tabular}{@{}>{\raggedright\arraybackslash}p{3.2cm}>{\raggedright\arraybackslash}p{6.2cm}>{\raggedright\arraybackslash}p{4.5cm}>{\raggedright\arraybackslash}p{4.5cm}@{}}
\toprule

\textbf{Axe d'action} &
\textbf{Proposition} &
\textbf{Objectif} &
\textbf{Intervenants} \\

\midrule

Sessions d'information &
Présenter en conseils municipaux les dispositifs de plans communaux de sauvegarde (\gls{pcs}). &
Renforcer la connaissance des dispositifs de gestion de crise au niveau local. &
Élus locaux \\

\addlinespace

Sessions d'information &
Mettre en place des actions de sensibilisation à destination des conseils municipaux des enfants, fondées sur une approche pédagogique adaptée à leur tranche d'âge. &
Développer une culture de résilience dès le plus jeune âge. &
Élus locaux, enseignants, intervenants spécialisés \\

\addlinespace

Sessions d'information &
Organiser des soirées thématiques à destination des élus locaux (maires notamment) et des sessions d'information pilotées par les préfets à destination des élus locaux, portant sur les dispositifs de gestion de crise. &
Renforcer la coordination et l'appropriation des dispositifs de crise au niveau local. &
Préfets \\

\addlinespace

Partenariats académiques et médiatiques, appuyés par l'IHEDN &
Impliquer les trinômes académiques et les écoles de journalisme afin de renforcer la culture de sécurité et la lutte contre les infox. &
Développer une culture partagée de vigilance et d'analyse de l'information. &
Auditeurs IHEDN, universitaires, formateurs, professionnels des médias \\

\addlinespace

Modules de formation &
Collaborer avec le CLEMI et la plateforme VIGINUM pour élaborer des modules de formation relatifs à la gestion de crise et à la détection des infox. &
Structurer une offre de formation adaptée aux enjeux contemporains. &
CLEMI, VIGINUM, experts en cybersécurité et information stratégique \\

\addlinespace

Intégration aux parcours IHEDN &
Intégrer les modules de formation développés aux parcours de formation des auditeurs de l'IHEDN. &
Assurer une montée en compétence progressive et homogène des auditeurs. &
Corps enseignant IHEDN, intervenants institutionnels \\

\bottomrule
\end{tabular}
}

\caption[Propositions d'actions de sensibilisation et de formation]{Propositions d'actions de sensibilisation et de formation contribuant à la préparation et à la résilience face aux crises majeures.}
\label{tab:sensibilisation_resilience}

\end{table}

\subsection{Pendant la crise : une place limitée}

Lorsque la crise est engagée, le cadre d'action se resserre rapidement. La gestion de crise obéit à une organisation strictement hiérarchisée, au sein de laquelle les auditeurs peuvent intervenir principalement au titre de leurs fonctions professionnelles ou de leur engagement opérationnel.

Lorsque l'événement est en cours, la contribution directe de l'association en tant que telle demeure limitée, l'expérience de ses membres s'exerçant principalement dans le cadre de leurs responsabilités propres.

\subsection{Après la crise : analyse et retour d'expérience}

En aval, les membres de l'association peuvent apporter une valeur ajoutée. Leur expérience et la diversité de leurs parcours peuvent nourrir l'analyse des crises, la participation aux retours d'expérience et l'élaboration de recommandations.

\begin{proposition}{Intégrer l'\ARSeize{} dans les démarches de retour d'expérience}
Le comité considère que l'\ARSeize{} pourrait contribuer aux démarches d'analyse post-crise et de retour d'expérience, en s'appuyant sur la diversité des expériences professionnelles et sectorielles de ses membres.
\end{proposition}

Cette lecture temporelle confirme que l'enjeu n'est pas seulement d'affirmer une disponibilité de principe. Il s'agit d'abord de mieux connaître les profils, les compétences et les contraintes des membres de l'\ARSeize, afin d'identifier les contributions possibles et les conditions dans lesquelles elles pourraient s'organiser.
